\subsubsection{Evaluierungsergebnisse der geometrischen Merkmale}

Nach der Evaluierung der reinen Lokalisierungsgenauigkeit schließt sich die Untersuchung der extrahierten geometrischen Parameter an. Die aus den 2D-Kronenpolygonen abgeleiteten Werte für die Baumhöhe und den Kronendurchmesser werden hierzu direkt den quantitativen Einträgen im städtischen Baumkataster gegenübergestellt. Die Ergebnisse sind in Tabelle \ref{tab:evaluierung_geom} dargestellt.

\begin{table}[htbp]
\centering
\caption{Ergebnisse der geometrischen Evaluierung in den Testgebieten Gohliser Straße und Mariannenpark}
\label{tab:evaluierung_geom}
\begin{tabular}{l c c c c}
\toprule
\textbf{} & \multicolumn{2}{c}{\textbf{Baumhöhe [m]}} & \multicolumn{2}{c}{\textbf{Kronendurchmesser [m]}} \\
\cmidrule(r){2-3} \cmidrule(l){4-5}
& Gohliser Str. & Mariannenpark & Gohliser Str. & Mariannenpark \\
\midrule
\textbf{Stichprobe (n)} & 44 & 306 & 44 & 306 \\
\addlinespace
\textbf{Korrelation (Pearson r)} & 0,82 & 0,72 & 0,71 & 0,29 \\
\addlinespace
\textbf{RMSE [m]} & 2,23 & 3,06 & 2,93 & 4,15 \\
\addlinespace
\textbf{Bias [m]} & -0,8 & -0,35 & -1,91 & -1,66 \\
\bottomrule
\end{tabular}
\end{table}

Die statistische Auswertung der Straßenbäume belegt eine starke lineare Korrelation der berechneten Baumhöhen mit den Katasterdaten bei moderater Abweichung mit einem RMSE von 2,23 m und einer geringen systematischen Unterschätzung (Bias = -0,80 m). Auch die Kronendurchmesser weisen im lichten Straßenraum mit r = 0,71 einen soliden Zusammenhang auf, wenngleich der negative Bias von -1,91 m eine engere geometrische Erfassung durch den Algorithmus anzeigt.

Im Mariannenpark bleibt die Korrelation der Baumhöhen mit r = 0,72 auf einem akzeptablen Niveau, während die Kronendurchmesser mit r = 0,29 und einem hohen RMSE von 4,15 m eine deutliche methodische Schwäche offenbaren. Die Kronenausdehnung wird mit einem Bias von -1,66 m auch in diesem dichten Bestand systematisch unterschätzt.

Eine primäre Ursache für die geometrischen Abweichungen liegt in den datenimmanenten und zeitlichen Differenzen zwischen den Datensätzen. Zwischen der LiDAR-Erfassung und dem Datenstand des Kataster besteht ein dreijähriger Zeitverzug, in dem natürliches Wachstum sowie städtische Baumpflegemaßnahmen die Kronengeometrie verändert haben. Zudem beruhen die Höhen- und Kronenmaße im amtlichen Katasters auf visuellen Schätzungen im Rahmen regulärer Baumkontrollen, wohingegen der LiDAR-Workflow die physische Dimension sensorisch direkt erfasst. Neben diesen Datendifferenzen führt die Bestandsdichte zu einer methodischen Fehlerfortpflanzung von der Positions- auf die Attributebene. Im dichten Kronendach des Mariannenparks treten vermehrt algorithmische Fehlzuordnungen auf. Verknüpft die Zuordnung einen detektierten Baum fälschlicherweise mit einer benachbarten Katasterreferenz einer anderen Größenklasse, verzerrt dieser unpassende Vergleich die Korrelation und den RMSE drastisch.

Insgesamt verdeutlichen die statistischen Kennzahlen die spezifischen Einsatzgrenzen des entwickelten Workflows. Während die starke Korrelation der Straßenbäume die hohe Eignung des Verfahrens für die automatisierte Inventarisierung isolierter Bäume belegt, zeigen die schwachen Werte des Parkbereichs die methodischen Grenzen der Segmentierung auf.